% §6.5 Confounder sensitivity, 9-city panel (Boston/NYC/SF analyzed
% separately in the original 3-city paper version).
% Numbers locked against results/replications/confounder_sensitivity_12.csv.

\subsection{Sensitivity to unobserved confounding and confounder-block ablation}
\label{sec:sensitivity-12}

The Shen-Ross DML estimates in Section~\ref{sec:shen-12city} rest on
the conditional ignorability assumption that, after partialling out
the structured property covariates and the location indicators we
observe, residual variation in the listing language is exogenous to
residual variation in the log sale price. This assumption is not
directly testable. We follow \citet{cinelli2020making} and bound the
strength of any unobserved confounder that would be needed to
overturn the per-market findings, and we run a leave-one-block-out
ablation that identifies which of the observed covariate groups is
doing the heavy lifting in each market.

\paragraph{Robustness Value.} The Robustness Value (RV) of
\citeauthor{cinelli2020making} quantifies the minimum strength of an
unobserved confounder, measured as its partial $R^2$ on both the
outcome and the treatment, that would suffice to overturn the
estimate.  Higher RV means more robust to unobserved confounding.
For each of the nine markets in which we have the post-deduplication
sensitivity machinery, Table~\ref{tab:confounder-sensitivity-12}
reports the RV alongside the DML point estimate and the Cinelli-Hazlett
$f^2$ measure of the strongest observed covariate's partial $R^2$.

Three markets exhibit RV values above $0.10$, indicating that an
unobserved confounder would need to explain more than $10\%$ of
residual variance in both the listing-uniqueness score and the log
sale price to overturn the estimate. Dallas has the highest RV at
$0.158$, followed by Phoenix at $0.110$ and Seattle at $0.100$.
The remaining six markets have RV values below $0.04$, indicating
that comparatively modest unobserved confounding could overturn the
per-market estimate. The pattern is consistent with the Section
\ref{sec:shen-12city} reading that Dallas is the second market in
which we can identify a Shen-Ross effect at the $95\%$ level: not
only is the point estimate large in absolute terms, the inference is
substantially more robust to confounders than in the markets where
the point estimate clusters near zero.

The $f^2$ values, which give the strength of the most influential
\emph{observed} covariate in each market on a comparable
partial-$R^2$ scale, sit between $1\text{e-}5$ and $0.030$ across the
nine markets. These observed-covariate benchmarks are between two
and four orders of magnitude below the RV thresholds in the
high-robustness markets, which means an unobserved confounder would
need to be substantially stronger than any covariate we observe in
order to flip the inference. This is reassuring in the high-RV
markets and uninformative in the low-RV markets, where neither the
observed nor the hypothetical-unobserved benchmarks are large.

\paragraph{Leave-one-block-out ablation.}
We partition the observed covariates into five blocks: census
demographics, crime and safety, points-of-interest amenities, spatial
controls (lat-lon and ZIP indicators), and micro-geographic features
(within-block parcel attributes). For each of the nine markets we
drop each block in turn and recompute the DML estimate, reporting the
maximum absolute change in $\hat\theta$ across the five
ablations. The block that produces the largest change is the
operative confounder block in that market.

The leave-one-block-out structure reveals striking cross-city
heterogeneity in which covariate block matters most.  In Seattle,
dropping the crime block moves the DML estimate from $+0.030$ to
$+0.511$, a change of $+0.48$ in $\hat\theta$ at the
per-$\sigma$ scale, by far the largest single-block sensitivity in
the panel.  This is consistent with the well-documented role of
crime indicators in Seattle's listing language and structured
price-formation process. In Dallas the operative block is census
demographics, dropping which moves $\hat\theta$ by $+0.015$.
Portland, Phoenix, Denver, and Washington DC are most sensitive to
the amenities block, with maximum block changes between $+0.004$ and
$-0.013$.  Atlanta and Chicago are most sensitive to spatial
controls (with changes of $+0.011$ and $+0.003$ respectively), and
Philadelphia is most sensitive to census demographics with a change
of $-0.011$.

\paragraph{Substantive interpretation.} Two readings of the ablation
pattern are worth distinguishing. The first attributes the
heterogeneity to local economic structure. Markets in which crime is
a salient price-relevant attribute (Seattle being the clearest case)
show the largest single-block sensitivity to dropping crime
indicators; markets in which amenity density is doing the relevant
price-formation work (Portland, Phoenix, Denver, DC) show the largest
sensitivity to dropping amenity controls. Under this reading the
appropriate confounder set for a per-market DML deployment is not a
single national specification but a market-specific specification
that prioritizes the locally-operative confounder block.

The second reading attributes the heterogeneity to noise. With $n =
290$--$334$ listings per market and $24$--$35$ covariates per block
within the partial-out regression, the relative contributions of the
five blocks are estimated with substantial noise, and the
``most-influential block'' identification may shift if we re-sample
the per-market panel. We cannot adjudicate cleanly between the two
readings at our current sample size. Both pre-registered guidance
and the meta-regression in Section~\ref{sec:meta-regression-12} would
favor the first reading at scale, and we report the per-market
ablations as suggestive rather than definitive.

\paragraph{Limitations.} The sensitivity analysis above is computed
on the nine cities of the post-deduplication panel. Boston, New York,
and San Francisco were analyzed under a different sensitivity
protocol in the three-city pilot reported in the original draft of
this manuscript and are not directly comparable to the nine-city
table above; we report the original three-city sensitivity in
Appendix~\ref{app:original-sensitivity} and note that a unified
twelve-city sensitivity protocol is a candidate for future work.
